\Needspace{12\baselineskip}
\fancyhead[L]{\small FIRST COMET OF 1805.}
\begin{center}
FIRST COMET OF THE YEAR 1805.

\rule{0.6\linewidth}{0.4pt}

Monatliche Correspondenz. Vol. XIII. pp. 79–83. January 1806.

\rule{0.6\linewidth}{0.4pt}

—\quad—\quad—
\end{center}

Dr. \textsc{Gauss} writes to us from Brunswick on 5 December that he has tried, without success, to fit these elements of \textsc{Bessel} to the observations from Frankfurt and Marseille. These observations, however, seem to be quite imprecise. “Since these elements,” writes Dr. \textsc{Gauss}, “differ by several degrees from \textsc{Huth}’s latest observations, I tried to improve them using those observations, even though \textsc{Huth}’s figures, given only in multiples of five minutes, seem to be mere estimates by eye. The attempt failed. The observations by \textsc{Thulis} that you were kind enough to send me (November issue, 1805, \href{https://www.digitale-sammlungen.de/en/view/bsb10538604?page=518}{p. 502}) do not seem particularly good either. At least I could obtain no satisfactory result from them. The comparison with \textsc{Olbers}’s observations already shows this. When reduced to the same time, his observations differ from those of \textsc{Thulis} on 29 October by six minutes in right ascension and nine minutes in declination, and on 31 October by one minute in right ascension and twelve minutes in declination. On 9 November, however, \textsc{Bessel}’s elements differ from \textsc{Thulis}’s observations by 71 minutes in longitude, which is more than can reasonably be allowed for observational errors. I therefore tried to correct the error, but I was still unable to avoid an error of almost 29 minutes in latitude. Here are the elements

\begin{center}
\begingroup\fontsize{8}{10}\selectfont\setlength{\tabcolsep}{3pt}\renewcommand{\arraystretch}{1.1}
\begin{tabular}{lr}
Passage through perihelion in November . . . . & $17.746$ in Seeberg. \\
Longitude of the ascending node . . . . . . . . & $340^\circ\ 11'$ \\
Inclination of the orbit . . . . . . . . . . . . . & $17\ 34$ \\
Longitude of the point nearest the Sun . . . . . . . . . . . . & $157\ 17$ \\
Logarithm of the least distance . . . . . . . & $9.53969$ \\
Motion . . . . . . . . . . . . . . . . & prograde \\
\end{tabular}\endgroup
\end{center}

“These results, however, warrant little confidence either.”
\SourceLeaf{271}
\par\indent We have mentioned all these circumstances solely to spare others working on the calculation fruitless effort should they try to bring the Frankfurt and Marseille observations into any kind of agreement. Since \textsc{Bessel}’s orbit represents \textsc{Olbers}’s observations up to 13 November so satisfactorily, these elements may be regarded as accurately determined. If we obtain better observations covering a longer period in the future, the elements may be expected to require only very slight improvements.

\hfill[\textsc{von Zach.}]

